\section{Introduction}
\label{sec:intro}

Autoregressive transformer inference caches the key and value projections of every past
token so that each new token attends over history without recomputation. This KV cache is
linear in sequence length, in the number of layers, and in the number of key/value heads.
For a model the size of Qwen2.5-1.5B the per-token KV state is small, but it accumulates:
at $64$k tokens the full-precision cache alone is larger than the model weights, and on a
shared $8$\,GB device it is the first resource to run out. The practical question this work
addresses is narrow and concrete: \emph{what does it take to run a long-context workload on
the hardware people actually own}---a laptop with unified memory and no discrete GPU---
rather than on a server with hundreds of gigabytes of HBM?

The standard answers each give up something. \emph{Quantizing} the cache to $4$ or $2$ bits
shrinks it by a constant factor but still grows linearly and eventually loses fidelity on
the tokens that matter. \emph{Evicting} tokens (keeping only ``heavy hitters'' or a sliding
window) bounds the cache but discards information that a later query may need.
\emph{Low-rank latent} caches change the model architecture and require retraining. We pursue
a different point in the design space, motivated directly by the structure of the cache: KV
vectors within a local block are highly correlated, so each block is well approximated by a
single representative token plus a low-rank correction. This is the central idea of
\diffkv{}: store an \emph{anchor} and a \emph{differential} low-rank term per block, and
never materialize the dense cache again---not even at attention time.

\paragraph{What \diffkv{} is.}
\diffkv{} is a runtime, not a model change. It monkey-patches the attention of a stock
quantized model and routes the two phases of inference differently. During \emph{prefill},
it runs exact attention while streaming the produced KV into a compressed store: every
$256$-token block is reduced to an anchor token plus a rank-$16$ joint-SVD factorization of
the per-token deltas (\S\ref{sec:compression}). During \emph{decode}, a single fused,
\texttt{mx.compile}-d kernel (\S\ref{sec:decode}) scores the incoming query against all
compressed blocks \emph{in low-rank coordinates}---projecting the query through each block's
$r$-dimensional basis rather than reconstructing $B{\times}d$ keys---and against a small
uncompressed recency window, then merges the two with a flash-style log-sum-exp. The
compressed pool is pre-allocated and fixed in size, so the KV state is bounded regardless of
how long the context grows (\S\ref{sec:runtime}).

\paragraph{What we measure, and how honestly.}
A central methodological commitment of this report is to attribute every number to the code
path that actually produced it. In the course of this work we found that an earlier
benchmark's headline throughput came from an \emph{exact} full-KV decode path, not from the
compressed kernel; conflating the two would credit the compression with speed and accuracy
it does not provide. We therefore report the compressed sparse decode as the primary result
and keep the exact decode as an explicit \emph{upper-bound ablation} on the same store
(\S\ref{sec:eval}). Likewise, because the global allocator peak is dominated by weights and
the transient prefill activations, we make the bounded \emph{KV-state footprint} the
headline memory metric---this is what the compression actually controls---and report the
allocator peak alongside it.

\paragraph{Contributions.}
\begin{enumerate}[leftmargin=1.4em,itemsep=2pt]
  \item \textbf{A differential block representation for KV} (anchor + rank-$r$ joint-SVD
  delta over normalized per-token residuals) that reduces KV state by $\approx\!10\times$
  per block at fp16, with a closed-form storage budget (\S\ref{sec:compression}).
  \item \textbf{A bounded-memory runtime} that streams compression during a chunked,
  numerically exact prefill, releases the prefill activations at the prefill$\to$decode
  boundary, and serves decode entirely from a fixed-size compressed pool plus a recency
  window (\S\ref{sec:runtime}).
  \item \textbf{A fused low-rank decode-attention kernel} that computes attention scores and
  outputs from the compressed factors without decompression, and merges the compressed and
  dense branches stably (\S\ref{sec:decode}).
  \item \textbf{An honest, reproducible evaluation} on an $8.6$\,GB Apple~M3 that separates
  the compressed path from an exact-decode upper bound, quantifies the memory/context-reach
  benefit and the throughput and fidelity costs, and is structured so a CUDA/Triton
  fused-decode path can be slotted in later (\S\ref{sec:eval}--\S\ref{sec:analysis}).
\end{enumerate}

\noindent
We deliberately scope this report to the active (MLX) runtime. A separate C++/GGML
reimplementation exists in the repository and is under active development; it is not the
subject of this report and none of its numbers appear here.
